% LSPV-5 EXPERIMENT::COVERBREAK1::2026-10-02T03:44:00-07:00
\documentclass[11pt]{article}\usepackage[margin=.88in]{geometry}\usepackage{amsmath,amssymb}
\begin{document}
\section*{COVERBREAK1: adversarial attack on the small-cover/maximal-exterior pattern}
\textbf{STATUS: COMPUTATIONAL / OPEN.} P versus NP remains open. The declared SU05/Hadamard finite control model is tested; no author-unique vMS binary sampler or asymptotic theorem is claimed.

\paragraph{Attack design.}
Separate \(m\) from \(h\); include the corrected RANKTRANS deficient ranges as positive controls; exhaust all nonconstant words for \((p,h)=(11,5),(13,7)\); and sample sparse, alternating and pseudorandom words at larger field ratios \((31,5)\) and \((127,5)\). Optimize lexicographically for large greedy cover, small exterior dimension \(d\), then small \(R\). Exact \(\tau_L\) is recomputed for the strongest hostile family.

\paragraph{Hostile family found.}
At \((p,h)=(13,7)\), Boolean grid word
\[
w=0000110
\]
gives, for every tested \(5\le m\le9\),
\[
\boxed{R=m,\qquad d=m,\qquad \tau_L=\lceil m/2\rceil.}
\]
The exact finite data are
\[
\begin{array}{c|ccccc}
m&5&6&7&8&9\\\hline
d&5&6&7&8&9\\
\tau_L&3&3&4&4&5\\
(d_i)&(2,2,1)&(2,2,2)&(2,2,2,1)&(2,2,2,2)&(2,2,2,2,1)
\end{array}
\]
and every legal slice in this family has rank at most two. Thus the hostile mechanism is genuinely realizable in the declared finite model: linear total span can coexist with linear exterior span and a linear-size minimum slice cover. No extrapolation beyond \(m\le9\) is made.

For this family DIMSUM2 yields only
\[
\rho(H)\ge1+2m,
\]
while SLICECOVER1 yields a comparable linear lower bound. Hence neither current structural inequality falsifies an \(O(m)\) contraction here. MATCIRC1 predicts abundant sparse common quadratic circuits once the exterior generator is acquired.

\paragraph{Positive control.}
The battery also reproduces the corrected RANKTRANS phenomenon. For \(p=11,h=5\), exterior deficiencies begin in the exhaustive tested family at \(m=5\); for \(p=13,h=7\), deficiencies are present throughout the relevant tested widths. Therefore the search is capable of detecting nonmaximal exterior structure and is not mechanically biased toward saturation.

\paragraph{Large-field side.}
For \(p=31,h=5\), tested \(4\le m\le10\), and \(p=127,h=5\), tested \(4\le m\le12\), every adversarial/random word tested had
\[
R=m,\qquad d=\binom m2,
\]
with greedy cover at most two. These larger-field tests were sampled, not exhaustive. They strengthen the finite transition hypothesis but do not prove any threshold in \(p/h\), \(p,h,m\), or the actual vMS regime.

\paragraph{Interpretation.}
COVERBAT1's maximal-rank pattern is not universal. The most useful new fact is the explicit finite hostile family: any future proof must use the large source-parameter separation or another source-specific property; nonconstancy and cyclic-window structure alone are insufficient. Conversely, the disappearance of deficiencies in the tested larger-field controls suggests attacking a quantitative large-field shortening/innovation lemma rather than a universal lemma.

\paragraph{Next target.}
Search for a theorem of the following guarded form: under the actual source inequalities (eventually \(h=m^{100}\), \(\Delta=hr\), \(p\ge\Delta^{100}\)), classify exceptional Boolean-grid leaves for which \(\dim U=O(m)\), or prove \(\dim U=\omega(m)\) outside that class. First extend adversarial computation with \(p/h\) increasing while preserving \(h>m\), because the current large-field controls use \(h=5\) and eventually \(m>h\), unlike vMS.

\paragraph{Artifact.}
Raw result: \texttt{results/raw/product\_span/2026-10-02\_LSPV5\_COVERBREAK1.json}. Reproducible core routines are in \texttt{experiments/product\_span/slice\_cover\_battery.py}; this run uses the same exact GF(2) span/exterior/cover calculations.
\end{document}
% END_OF_VERIFIED_EXPERIMENT::LSPV-5::COVERBREAK1